% --- ページ番号の設定 (1/3, 2/3形式) ---
\usepackage{lastpage}
\usepackage{fancyhdr}
\pagestyle{fancy}
\fancyhf{}
\cfoot{\thepage\ / \pageref{LastPage}}
\renewcommand{\headrulewidth}{0pt}

\begin{document}
\maketitle

\section*{[課題1-1] LLMの活用：検索との比較と価値創造}

\subsection*{1. 検索エンジンとLLMの比較分析}
【調べた内容：例「大規模言語モデルにおける感情理解の現状」など】
\begin{itemize}[nosep]
    \item \textbf{検索の結果：} 特定の論文や技術ブログへのリンクが得られるが、断片的な情報の統合は人間に委ねられる。
    \item \textbf{LLMの結果：} 概念の要約、歴史的背景、現在の課題が体系的に整理されて出力される。
    \item \textbf{意味合い：} 検索は「点（情報の所在）」を探すのに適し、LLMは「面（構造の理解）」を構築するのに適している。
\end{itemize}

\subsection*{2. 「役に立つ」プロンプトの共有}
【お題：例「Polaris開発に向けた倫理的ロードマップの策定」など】
\point{選択した問い}{「AIが人間の複雑な感情を模倣する際、法的に保護されるべきプライバシーの境界線はどこにあるか？」}
\point{有用性の理由}{法的・倫理的観点を多角的に提示し、開発者が盲点になりやすい「感情データの所有権」についての議論を深めてくれたため。}

% ここにキャプチャ画像を貼る（1枚に収めること）
\begin{figure}[H]
    \centering
    \includegraphics[width=0.75\textwidth]{../Assets/llm_capture.png}
    \caption{LLMとの対話ログ（思考の深化プロセス）}
\end{figure}

\clearpage % 2ページ目（課題1-2）へ強制改行
% ここから先ほどまで作っていた「分析」の内容を続ける